\section{Regulatory, Ethical, and Oversight Considerations}

The oversight architecture for this combined IND/IDE first-in-human study layers a
conventional human-subject-protection structure with the additional governance
that an autonomy-bearing surgical system requires. The Sponsor-Investigator,
ChemicalQDevice, answers to a reviewing Institutional Review Board (IRB) and to an
independent data and safety monitoring board (DSMB), convenes a Physical AI Safety
Review Committee on a $\leq 90$-day cadence, designates an Independent Safety
Monitor (ISM) for each procedure, and delegates day-to-day monitoring to a
robotics- and artificial-intelligence-competent clinical research organization
(CRO) over the site team. Figure~16 presents the governance and safety-oversight
structure.

\begin{mermaidfig}
\node[mmgoal](si)  at (4.2,0)     {Sponsor-Investigator\\ChemicalQDevice};
\node[mmstep](irb) at (0,-2.4)    {Reviewing IRB\\SR device + IND};
\node[mmstep](dsmb)at (4.2,-2.4)  {Data and Safety\\Monitoring Board};
\node[mmstep](pais)at (8.4,-2.4)  {Physical AI Safety\\Review Committee\\$\leq 90$-day cadence};
\node[mmin]  (cro) at (-2,-4.8)    {CRO / clinical\\monitor (robotics\\+ AI competent)};
\node[mmstep](ism) at (8.4,-4.8)  {Independent Safety\\Monitor (per\\procedure)};
\node[mmin]  (site)at (4.2,-6.5)  {Site team: surgeon,\\operator, backup\\operator};
\draw[mmedge] (si) -- (irb);
\draw[mmedge] (si) -- (dsmb);
\draw[mmedge] (si) -- (pais);
\draw[mmedge] (si.west) -| (cro.north);
\draw[mmedge] (cro) |- (site.west);
\draw[mmedged] (dsmb) -- node[right,font=\scriptsize]{stop/continue} (site);
\draw[mmedged] (pais.south) -- node[right,font=\scriptsize,pos=0.5]{USL changes} (ism.north);
\draw[mmedged] (ism) |- (site.east);
\end{mermaidfig}
\vspace{-0.6cm}
\figcaption{Figure 16. Study governance and safety oversight. The
Sponsor-Investigator answers to the reviewing IRB\\and the DSMB, convenes a
Physical AI Safety Review Committee on a $\leq 90$-day cadence, designates\\an Independent Safety Monitor for each procedure, and delegates monitoring to a
robotics- and\\AI-competent CRO over the site team (surgeon, operator, and a
qualified backup operator).}

\subsection{Informed Consent Process}

Informed consent is obtained by a qualified member of the research team, who is
not the treating surgeon, before any study-specific procedure, in language the
participant can understand and with adequate time and opportunity to ask
questions, consistent with 21 CFR part 50 \cite{ecfr2026part50} and ICH E6(R3).
The consent discussion covers the diagnosis, the alternatives, and the prognosis;
the perioperative daraxonrasib regimen and its risks; and, in addition to the
standard elements, an explicit and distinct disclosure of the on-premises
LLM-directed robotic system, its role and continuous human oversight, and its
system-specific risks, as required for an autonomy-bearing investigational device.
The consent form addresses the therapeutic misconception directly and explains the
staged-autonomy model under which the system operates only at the
clinician-controlled and supervised semiautonomous levels.

A defining feature of this trial class is the documented Physical AI opt-out.
Under 21 CFR $\S$312.60(f), as adapted for Physical AI investigations
\cite{kawchak2026adapt312}, the participant is offered, and may exercise, a
documented right to request administration without the Physical AI advisory layer
where clinically feasible, without penalty and without affecting eligibility for
the remainder of the study or for clinical care. The participant's election or
declination of the opt-out is recorded in the source document and the case report
form (CRF). Assent is obtained where applicable, although the adult-only
eligibility makes pediatric assent inapplicable in practice; a legally authorized
representative process and the use of qualified interpreters and translated
materials for participants with limited English proficiency are provided so that
consent is meaningful for every eligible individual. The right to withdraw at any
time without penalty is stated at consent and reaffirmed at each contact.
Figure~17 shows the consent process and the opt-out.

\begin{mermaidfig}
\node[mmin]  (intro)at (-0.2,0)      {Investigator presents\\study: disease,\\alternatives, prognosis};
\node[mmstep](pai)  at (4.4,0)    {Explain Physical AI\\system: role, oversight,\\risks (\S312.60(f))};
\node[mmdec] (opt)  at (9.3,0)    {Participant\\choice};
\node[mmgoal](yes)  at (9.3,-2.6) {Consents to Physical AI\\signed ICF + assent\\if applicable};
\node[mmstep](no)   at (4.4,-2.6) {Requests non-Physical AI\\administration if\\clinically feasible};
\node[mmgoal](enr)  at (4.4,-4.9) {Enroll and document\\consent in source + CRF};
\node[mmin]  (wd)   at (-0.2,-4.9)   {Right to withdraw at\\any time without penalty};
\draw[mmedge] (intro) -- (pai);
\draw[mmedge] (pai) -- (opt);
\draw[mmedge] (opt) -- node[right,font=\scriptsize]{accept} (yes);
\draw[mmedge] (opt.south west) -- node[above,font=\scriptsize, xshift=-0.42cm, yshift=0.25cm]{decline AI} (no.east);
\draw[mmedge] (yes) |- (enr.east);
\draw[mmedge] (no) -- (enr);
\draw[mmedged] (enr) -- (wd);
\end{mermaidfig}
\vspace{-0.65cm}
\figcaption{Figure 17. Informed-consent process with the Physical AI opt-out.
Consent discloses the Physical AI system's role, continuous human oversight, and
system-specific risks, and offers the participant a documented right to request
non-Physical AI administration where clinically feasible (21 CFR $\S$312.60(f)),
with standard right to withdraw.}

\subsection{Study Discontinuation and Closure}

The study may be discontinued or closed by the Sponsor-Investigator, by the IRB,
or by the FDA. Grounds for closure include unacceptable safety findings (including
the prespecified halt rules of \S9.4.6), a clinical hold, a determination that the
study cannot be conducted in compliance, or completion of the planned cohorts. The
FDA may impose a clinical hold on the drug or device component under the
provisions for terminating or suspending an investigation; in addition, the
Physical AI overlay establishes specific clinical-hold grounds for the
autonomy-bearing system, including a breach of a hard cyber-physical safety limit,
a failure of the verification-before-generation gate surface, or a loss of audit
integrity. On any hold or closure, enrollment and patient-facing robotic activity
stop immediately; affected participants are transitioned to appropriate clinical
care; and a controlled decommissioning and data-preservation procedure is executed,
under which the robotic system is locked out of patient contact, the deterministic
configuration and repository commit are recorded, and the complete hash-chained
command and telemetry record is sealed and retained for the regulatory retention
period (\S10.9). The IRB, the FDA, and the DSMB are notified, and a final report is
prepared.

\subsection{Confidentiality and Privacy}

Participant confidentiality is protected throughout. Data released for analysis,
deposition, or publication are de-identified by the HIPAA Safe Harbor method, with
removal of all eighteen categories of protected health information identifiers
(names; geographic subdivisions smaller than a state; all date elements other than
year that are directly related to an individual, and ages over 89; telephone and
fax numbers; email addresses; Social Security numbers; medical record numbers;
health plan beneficiary numbers; account numbers; certificate or license numbers;
vehicle identifiers; device identifiers and serial numbers; web URLs; IP
addresses; biometric identifiers; full-face photographs and comparable images;
and any other unique identifying number, characteristic, or code). Access to
identifiable records is governed by a deny-by-default authorization model, in which
no role can read, write, or export a record except where an access right has been
explicitly granted, and every access, change, and export is written to a
hash-chained audit trail in which each entry cryptographically incorporates the
hash of the prior entry so that the record is tamper-evident and any alteration is
detectable. All electronic records, signatures, and the cyber-physical telemetry
stream are maintained under 21 CFR part 11 \cite{ecfr2026part11}, with validated
systems, secure time-stamped audit trails that do not obscure prior entries,
role-based access controls, and electronic signatures bound to the records they
authenticate.

\subsection{Future Use of Stored Specimens and Data}

The cyber-physical record is retained as a tiered telemetry pyramid that balances
long-term scientific reuse against storage cost. At the apex are the compact,
human-readable summaries and the safety-event and audit logs retained in full and
indefinitely; the middle tiers hold compressed feature and event streams; and the
base, level zero (L0), holds the raw, losslessly reconstructable command and
sensor stream. The L0 raw tier is deposited on Zenodo under controlled terms and
is decompressible and reconstructable on FDA request as required for inspection
and review under 21 CFR $\S$312.57, so that any reported behavior can be replayed
deterministically from the seed and the repository commit. Any future secondary
use of stored de-identified data is limited to research consistent with the
consent and applicable regulation, is conducted on de-identified data only, and
preserves the audit chain. No biological specimens are banked beyond those required
for the clinical care and the protocol-specified assessments; tumor genotyping is
performed on standard-of-care material, and no specimen is retained for an
unspecified future purpose without appropriate consent.

\subsection{Key Roles and Study Governance}

The Sponsor-Investigator (ChemicalQDevice, Kevin Kawchak) holds the combined
IND/IDE and bears the sponsor and investigator responsibilities under 21 CFR parts
312 and 812, including protocol custody, safety reporting, oversight of the
investigation, and assurance of compliance. The site Principal Investigator
directs the conduct of the study at the clinical site and supervises the research
team. The operating surgeon performs the pancreaticoduodenectomy under the
staged-autonomy model; a primary operator drives the system, and a qualified
backup operator is present and ready to assume control or to invoke the manual
override at any time. The CRO provides robotics- and AI-competent clinical
monitoring. The DSMB, the Independent Safety Monitor, and the Physical AI Safety
Review Committee provide the independent oversight described in \S10.6. Roles,
delegated duties, and the corresponding qualifications are documented on a
delegation-of-authority log maintained by the site, and changes are recorded and
reported as required.

\subsection{Safety Oversight}

Independent safety oversight operates on three tiers. The DSMB, composed of
members independent of the conduct of the study with surgical, oncologic,
biostatistical, and device-safety expertise, reviews accumulating safety data at
each cohort boundary and at any triggered review, applies the prespecified halt
rules of \S9.4.6, and recommends continuation, modification, pause, or halt. The
Independent Safety Monitor is designated for each procedure and provides
real-time, case-level safety oversight independent of the operating team, with the
authority to call a stop. The Physical AI Safety Review Committee meets on a
$\leq 90$-day cadence (and ad hoc on any triggering event) to review the
autonomy-bearing behavior of the system: the telemetry and audit record, the
verification-before-generation gate-surface decisions, the Unified Safety Level rating, and any software change, and to recommend USL reassessment or
oversight changes to the Sponsor-Investigator. The three tiers are coordinated so
that no autonomy-bearing change and no escalation decision, proceeds without
appropriate independent review.

\subsection{Clinical Monitoring}

The study is monitored by a CRO whose monitors are competent in both clinical-trial
monitoring and the robotics and artificial-intelligence elements of this trial
class. Monitoring follows a risk-based plan and includes source-data verification
not only of conventional clinical source documents and the consent records but
also of the cyber-physical evidence: the device telemetry, the
verification-before-generation gate-surface logs, and the hash-chained audit trail,
which are verified against the recorded outcomes and the deterministic seed and
repository commit. Monitoring visits verify eligibility, consent (including the
documented Physical AI opt-out election), safety reporting, protocol adherence,
and data integrity, and findings are documented and resolved through to closure.

\subsection{Quality Assurance and Quality Control}

Quality assurance and quality control span the clinical and the cyber-physical
domains. Before every procedure, the pre-procedure safety matrix is verified, the
USL readiness rating is confirmed at or above the deployment threshold, and the
emergency-stop, positional, and force envelopes are checked. After any software
update, model change, or configuration change to the LLM-directed system, a
Phase 0 re-validation is performed against the independent simulation frameworks,
and the USL is reassessed, before any further patient-facing use; no update reaches
a participant without re-validation. Audits may be conducted by the
Sponsor-Investigator or a designee independent of the conduct of the study, and
corrective and preventive actions are documented. The deterministic, version-controlled
build ensures that the validated configuration in use is exactly the configuration
that was tested and recorded.

\subsection{Data Handling and Record Keeping}

Seven Physical AI record types are captured and retained for this trial class, in
addition to the conventional clinical trial records: (1) the LLM advisory record
(each prompt, context, and advisory output); (2) the robot command record (each
command issued to the arms); (3) the sensor and telemetry record (the
cyber-physical stream, tiered as in \S10.4); (4) the safety-limit-event record
(emergency-stop activations and latencies, positional and force-envelope
excursions, and vascular-exclusion gating events); (5) the
verification-before-generation gate-surface record (each accept, block, or
escalate decision); (6) the human-override and intervention record; and (7) the
hash-chained audit trail that binds the preceding six to a deterministic seed and a
repository commit. All records are retained for the period required by 21 CFR
$\S$312.57 (at least two years after a marketing application is approved for
indication, or, if no application is filed or approved, two years after the
investigation is discontinued and the FDA is notified), and the L0 raw tier is
preserved described in \S10.4 so that a requested record can be reconstructed.

\subsection{Protocol Deviations}

A protocol deviation is any departure from the IRB-approved protocol. Deviations
are identified, documented, assessed for their effect on participant safety and
data integrity, and reported to the IRB and the Sponsor-Investigator in accordance
with the IRB's policy and applicable regulation; deviations that affect safety or
the rights and welfare of participants, and any that recur, are escalated.
Deviations are distinguished from planned protocol amendments, which follow the
amendment process. Where a deviation is necessary to eliminate an apparent
immediate hazard to a participant, the change may be implemented without prior IRB
or FDA approval under 21 CFR $\S$312.30(h) (and the corresponding device
provision), with prompt notification and reporting afterward; for the
autonomy-bearing system, an immediate-hazard modification also triggers the
decommissioning and re-validation controls of \S10.2 and \S10.8 before any further
patient-facing use.

\subsection{Publication and Data Sharing Policy}

Results are disseminated regardless of outcome through dual deposition: the
analysis code, derived data sets, and protocol are published in a public GitHub
repository and archived on Zenodo with a citable DOI, under the Creative Commons
Attribution 4.0 International (CC BY 4.0) license. Reporting follows the TRIPOD+AI
statement for the model-assisted advisory component \cite{Collins2024tripodai} and
the CREMLS (Consensus Reporting for Evaluation of Machine Learning Systems)
guidance for the machine-learning system, in addition to the applicable
clinical-trial reporting standards. The deterministic seed and the repository
commit are published so that every reported result is reproducible, and the tiered
telemetry record is deposited as described in \S10.4. No identifiable participant
information is published; all shared data are de-identified by the HIPAA Safe
Harbor method (\S10.3).

\subsection{Conflict of Interest Policy}

Financial interests are disclosed and managed in accordance with 21 CFR part 54
(Financial Disclosure by Clinical Investigators). The Sponsor-Investigator is the
chief executive of ChemicalQDevice, which develops the Physical AI system under
study; this relationship is disclosed in full, and the independent oversight
structure of \S10.6 (the DSMB, the Independent Safety Monitor, and the Physical AI
Safety Review Committee) is established in part to manage it, so that safety and
escalation decisions are not made by a party with a financial interest in the
outcome alone. Investigator financial-disclosure certifications are collected
before participation and updated through one year after the study, and any
identified conflict is disclosed and managed before it can affect the conduct,
analysis, or reporting of the study.
